\documentclass[11pt]{article}
\usepackage[letterpaper,margin=.68in,headheight=16pt,headsep=.18in,footskip=.28in]{geometry}
\usepackage{amsmath,array,enumitem,fancyhdr,hyperref}
\pagestyle{fancy}\fancyhf{}\renewcommand{\headrulewidth}{.3pt}
\fancyhead[L]{\small Bellingham Math Circle / Week 11 / Return visits}
\fancyhead[R]{\small Adult guide}
\fancyfoot[L]{\small RV-11-FAC-v1 / Draft and unpiloted}\fancyfoot[R]{\small\thepage}
\setlength{\emergencystretch}{2em}
\setlength{\parindent}{0pt}\setlength{\parskip}{7pt}
\setlist[itemize]{leftmargin=1.3em,itemsep=3pt,topsep=4pt}
\hypersetup{hidelinks,pdftitle={Week 11 return-visit facilitator guide}}
\begin{document}
{\Large\bfseries Mathematical overview}\par\medskip
\textbf{Facts and limits.} On a finite undirected sink graph with finite nonnegative chip piles, each legal firing sends one chip along every incident edge. With every working vertex connected to an absorbing sink, every legal run continued until no move is available terminates, with the same finish and firing counts. Without a sink, termination may fail. If a start does admit legal stabilization, its complete legal runs still have unique finish and firing counts, even without a sink. On a closed triangle, total chips are conserved: a stable placement holds at most three. Even at that total, some starts stabilize and others cycle.

A one-chip addition to a stable state can start a cascade. On the four-cycle \(S-A-B-C-S\), all working thresholds are two. The largest avalanche from a stable start has four firings, not three: the middle vertex can fire twice. This maximum is specific to this board and the one-chip perturbation.

A finish does not determine its past. On the sink triangle each firing loses one chip to the sink. Finishing at \((1,1)\) after four firings fixes the initial total at six, but leaves three possible starts and eight legal histories. Formal balance alone does not guarantee a proposed word is legal.

\textbf{Children's destination.} Separate the reason a run stops from experiments where it happens to stop; notice propagation through occupied neighbors; reconstruct and check legal histories rather than assuming a unique past.

\textbf{Entry map.} Problem 1: K--1 with oral rules, small counting and an adult recording the run. Problem 2: grades 2--3 and up, distinguish a stable start from its one-chip perturbation. Problem 3: grades 4--5, small arithmetic and backwards reasoning; a counter reconstruction is sufficient, without algebra.

\textbf{What is new.} The base tests sink-board order independence, staged additions and repeated addition cycles. Problem 1 expands the sink-free cycle seed already suggested in the base adult guide's optional closed-board continuation (p. 12 as read October 4, 2026) into contrasting stabilization, cycle and capacity work. Problem 2 optimizes cascade length from stable starts. Problem 3 reuses the six-chip starts from base grades 2--3 Problem 3, but its new target is every legal firing history and the missing orders. They are three investigations within one chip-firing theme.
\newpage {\Large\bfseries Prepare a return visit}\par\medskip
\textbf{Status and use.} Three investigations extend one familiar theme. These companions are separate from the base packets and may support several visits. All pages are new and unpiloted. Printed labels describe prerequisites, not age restrictions. Give one page at a time; completing every page is not the goal.

\textbf{Materials and preparation.} For each working pair, 12 counters no wider than 20 mm, one board, pencil and spare paper/whiteboard. For the current five pairs plus upper-table third child, reserve 72 counters so each table can compare two starts. Printed working circles should be at least 42 mm and sinks at least 55 mm; stack piles rather than spread chips across lines. No exact fit or existing stock has been physically checked.

\textbf{Staffing.} Current context: eleven children, fixed KK11 / 3333 / 445 tables, one adult anchored to each table. Use pairs; the upper third child rotates as reader or checker. Routine adult reading, arrow drawing or recording can leave the mathematical decisions with children. If adults must choose every move, use the earlier entry rather than run the investigation for them.

\textbf{Short common launch.} Let children handle chips freely. On a sink board, place a small legal pile and ask a child to send one chip along every touching line. Show that extras remain and the sink never sends chips back. Write the fired circle’s letter after each move; a second letter is appended, not substituted. Then cover the sink label with C and show that the sink-free triangle is a separate board with a different rule, not a continuation of the old trial. A partner touches all recipient lines before the chips move.

\textbf{Suggested first route.} Start Problem 1 with three chips; offer the four-chip impossibility question after a real run. Start Problem 2 with an empty stable board and a single addition, then a board with some ones. Save Problem 3 for children who can restart a trial reliably and keep a four-letter move record. If they completed base grades 2--3 Problem 3, start from their known three starting states and investigate all legal histories rather than repeat that search.

\textbf{Flexible timing.} Allow 3--5 minutes of material play and 2--4 minutes for the common demonstration, then 15--25 minutes for one investigation and 5 minutes to share. A longer second visit can spend 30--45 minutes on the same investigation, including unfinished examples or its explanation. These are planning estimates, not observed timings. Do not require all three within one hour.

\textbf{Questions and hints.} Start with ``What can you try?'' and ``What would count as success?'' Check the actual rule before giving a strategy. Optional questions and hints are keyed below; use them only after a real attempt. A counter argument or drawing may be a complete explanation.

\textbf{Source and pilot record.} Mathematical examples and arguments here are original local follow-ups to the current base theme. Consulted teaching source: \emph{Math Circle by the Bay}, Preface pp. vii--x (local PDF pp. 8--11), on interaction, manipulatives, hints, explanations and themes spanning multiple years. Our exact launch, entry route and timing are design inferences. Year-1 \emph{Handouts 6}, Problems 6.2--6.6, provided a local example of connecting representations. Record actual problem, starting route, examples tried, what needed adult rescue, and what children want to resume. Distinguish observations from hypotheses. Counter fit, materials stock, procedures and classroom response have not been physically tested.
\newpage
\textbf{Problem 1: removing the escape.} On the closed triangle every vertex has degree two. From \((3,0,0)\), firing A gives \((1,1,1)\), which is stable. From \((2,1,0)\), the forced run is
\[
(2,1,0)\xrightarrow{A}(0,2,1)\xrightarrow{B}(1,0,2)\xrightarrow{C}(2,1,0).
\]
The same three legal moves can repeat forever. This is a certificate of nontermination, not just a long unfinished experiment.

With four chips, no stable state exists: each of three circles must hold at most one, while firing never changes the total. Therefore every continued legal run is infinite. There are only finitely many placements of four chips, so an earlier arrangement must repeat. For the printed start, AABC reaches (2,1,1), which was already reached after the first A; ABC can then repeat. With three chips the same capacity argument does not decide the outcome; the contrasting starts above are essential.

\textbf{If needed.} Keep a photograph or drawing of the start beside the board. Once the arrangement repeats, ask what the same next move will do. To show the capacity obstruction, put one chip in each circle and ask where the fourth could remain without a legal move.

\textbf{Return question.} Which three-chip starts stabilize? Up to permuting the circles, the only types are \((3,0,0)\), \((2,1,0)\), and \((1,1,1)\). The first reaches \((1,1,1)\), the second cycles, and the third is already stable. Children can give the proof with piles; ordered notation is adult support.

\textbf{Observe.} Can the child identify a repeated whole arrangement, or merely repeat a preferred move? Record that distinction before changing the activity.
\newpage
\textbf{Problem 2: the biggest one-chip avalanche.} Stable starts have zero or one chip at each of A, B and C. Add exactly one chip, then finish all legal sharing. The maximum number of firings is \textbf{four}, attained only by starting at \((1,1,1)\) and adding at B. Its two legal words are BACB and BCAB, each with counts \((1,2,1)\); the finish is \((1,0,1)\), and two chips have escaped.

Adding at A to the all-one start gives ABC and finish \((1,1,0)\). Adding at C gives CBA and finish \((0,1,1)\). Each has three firings. Every other stable start gives at most two. An addition to an empty circle gives none.

\textbf{Why the bound is complete.} There are eight stable starts and three addition sites, hence 24 trials. The independent check explores every legal order for each. A child can also reason locally: from an occupied endpoint, a wave can advance across occupied circles, but a missing one stops it. From B, each occupied neighbor can send one chip back; B fires again only when both neighbors initially held one. That is the all-one start, and after B's second firing both endpoints have just one chip.

\textbf{Light record.} Write the move word while moving, then count letters afterwards. A partner handles the counters or reads the next legal choice; do not require a concurrent tally.

\textbf{If needed.} Offer a start with a single empty neighbor. Ask where a second chip must come from before that neighbor can fire. Keep the starting board unchanged as a reference.

\textbf{Return question.} How does the maximum change on a longer path with a sink at both ends? Do not transfer the four-firing bound to a new board; this is a fresh conjecture to test.
\newpage
\textbf{Problem 3: all legal four-move histories.} The page uses the three known starts from base grades 2--3 Problem 3. The complete history answer is:
\begin{center}\begin{tabular}{c|l}
Start (A,B)&Legal four-letter histories\\\hline
(0,6)&BBAB, BBBA\\
(3,3)&ABAB, ABBA, BAAB, BABA\\
(6,0)&AAAB, AABA
\end{tabular}\end{center}
Every history ends \((1,1)\), with four chips in the initially empty sink.

\textbf{Completeness for the printed task.} For each supplied start, try all four-letter A/B words, rejecting a word at its first illegal firing. The table lists every surviving word that stops at the required finish. A pair can organize the attempts by their first letter, keeping just the word and final board.

\textbf{Readiness-dependent inverse continuation.} If the starts were not supplied, could they be recovered? Four moves lose four chips, so the start has six. Let A fire \(a\) times and B fire \(b\) times; \(a+b=4\). Reversing the changes gives starting piles \(1+2a-b\) and \(1+2b-a\). Counts \((0,4)\) and \((4,0)\) require a negative start, so only \((a,b)=(1,3),(2,2),(3,1)\) remain. These give the three starts in the table. Replay each multiset of letters, discarding a word immediately when its next firing is illegal. This leaves exactly the eight listed words.

A child need not use the formulas: reverse a firing by taking one chip from the other working circle and one from the sink and returning both to the named circle. Four reverse moves consume the four sink chips. A reverse move with an empty recipient is forbidden. Restore a copy after each branch; adult recording can support the tree of attempts.

\textbf{If needed.} Begin with a known forward four-move run, cover its start, then reconstruct it. The discovery task still concerns all possible pasts.

\textbf{Extension.} Change the known finish or run length. Keep the same exact information in every trial. A finish and a run length can determine the initial total without determining the initial arrangement.

\textbf{Checks.} Independent exhaustive firing checks are included in the source package. The algebra and reverse-counter explanation supply the reason for arbitrary order; code verifies this finite answer.
\end{document}
